\section{What TTA Failure Reveals About Physics-Structured Shifts}
\label{sec:rootcause}

\subsection{Defining Physics-Structured Shift}

We formalize the shift class that WiFi-TTA-Bench targets. Let $\mathbf{X} = G(\theta, \epsilon)$ be the data-generating process, where $\theta \in \Theta$ are physical parameters (path gains $h_l$, delays $\tau_l$, room geometry) and $\epsilon$ is observation noise. A \textbf{physics-structured shift} occurs when the environment changes $\theta \to \theta'$ while the governing equation $G$ remains invariant:

\begin{equation}
    \text{Source: } \mathbf{X}_s = G(\theta_s, \epsilon), \quad \text{Target: } \mathbf{X}_t = G(\theta_t, \epsilon), \quad G \text{ unchanged.}
\end{equation}

This differs from the standard covariate-shift assumption ($P(\mathbf{X})$ changes, $P(Y|\mathbf{X})$ preserved) and from corruption shift (additive perturbation to $\mathbf{X}$). In physics-structured shift, the mechanism generating features changes through unmeasured physical parameters, which can alter class-conditional structure. Our operational test for this paper is a cross-environment split in which room/session parameters change while the sensing task and label set are held fixed.

Table~\ref{tab:shift_props} quantifies this distinction:

\begin{table}[h]
\centering
\caption{Physics-structured shift is high-variance but structured: 231$\times$ more feature spread than synthetic shift, yet lower accuracy drop.}
\label{tab:shift_props}
\small
\begin{tabular}{lccc}
\toprule
\textbf{Property} & \textbf{Synthetic} & \textbf{Widar (real)} & \textbf{Ratio} \\
\midrule
Feature shift (L2) & 0.14 & 1.72 & 12$\times$ \\
Feature variance & 0.004 & 0.926 & 231$\times$ \\
Confidence std & 0.006 & 0.161 & 27$\times$ \\
Accuracy drop & 22.8~pp & 11.1~pp & 0.5$\times$ \\
\bottomrule
\end{tabular}
\end{table}

The paradox---high variance but lower accuracy drop---is consistent with physics-structured shift: features scatter widely in patterns governed by room geometry rather than by independent random corruptions.

\subsection{Why Entropy Minimization Can Fail Under Physics-Structured Shift}

Standard TTA is usually motivated by covariate shift: $P_t(\mathbf{X}) \neq P_s(\mathbf{X})$ but $P_t(Y|\mathbf{X}) = P_s(Y|\mathbf{X})$. Under physics-structured shift, the data-generating process $G(\theta, \epsilon)$ changes through $\theta$, which can alter both $P(\mathbf{X})$ and the effective class-conditional distributions. In that regime entropy minimisation reinforces the model's source-trained decision surface on target inputs even when that surface is misaligned; combined with overconfident source predictions (86--94\% confidence at 10\% accuracy), each entropy step can move predictions further from the target truth.

\subsection{Three Mechanisms of TTA Failure}

Evidence is consistent with three mechanisms, each with quantitative support. Hard physical-shift interpretations apply specifically to Widar\_BVP, whose environmental splits are verified rooms; NTU-Fi and SignFi use temporal/session splits and are treated as softer evidence.

\paragraph{Mechanism 1: Overconfident misclassification.} On Widar target rooms, aggregated across 15 runs, the source model's mean max-softmax confidence is 70--81\,\% across every class while per-class accuracy ranges from 24\,\% (draw-zigzag) to 57\,\% (push\&pull) (Appendix Table~\ref{tab:overconfidence}). Entropy-minimising adaptation reinforces these over-confident and often incorrect predictions, consistent with the 0.93--1.00 negative-adaptation rates we observe for entropy-family methods on Widar. The step-budget analysis (Appendix~\ref{app:step_budget}) shows monotonic degradation, consistent with each step pushing predictions further into incorrect modes.

\paragraph{Mechanism 2: Batch-statistic unreliability.} Faithful TENT on MLP+BN achieves $-$0.61~pp (Table~\ref{tab:faithful_tent}), which is consistent with BatchNorm plus BN-affine-only updates being insufficient on Widar\_BVP. BN-reset (forward pass only) is actively harmful ($-$1.11~pp), which is consistent with target running statistics diverging from source statistics at WiFi scale (400--3,000 samples per environment, 32--64 per batch).

\paragraph{Mechanism 3: Physics-parameter dependence.} The large swing (synthetic $+$7.4 to $+$8.0 $\to$ real $-$0.5~pp for the \texttt{physics\_tta} objective family) is consistent with physics-guided adaptation working when parameters $(h_l, \tau_l)$ are observable but failing without them, under our protocol. This motivates channel-estimation-integrated TTA without claiming that parameter observability is the only bottleneck.

\subsection{Label Smoothing: Source Calibration Changes TTA Outcomes}

Our label smoothing ablation (Table~\ref{tab:smoothing}) reveals that Mechanism~1 is partially addressable: smoothing $\alpha{=}0.1$ reduces \texttt{physics\_tta} damage from $-$0.20 to $-$0.10~pp and lowers negative adaptation rate from 0.60 to 0.47. The effect is modest but directionally consistent: \textbf{source training decisions propagate into deployment-time adaptation outcomes}.

\subsection{Research Agenda}

WiFi-TTA-Bench is designed to track three directions: (1)~channel-estimation-integrated TTA that consumes pilot-based $(h_l,\tau_l)$ estimates rather than unlabelled statistics alone; (2)~calibration-aware source training (label smoothing, mixup, temperature scaling, focal loss), guided by the exploratory smoothing effect ($-$0.20 $\to$ $-$0.10~pp); (3)~harm-aware adaptation protocols that target non-degradation rather than mean gain, given the 0.33--1.00 negative-adaptation variance we observe. The mean cross-\emph{dataset} method-rank Spearman $\rho{\approx}0.70$ (intro Finding 3) and a separate cross-\emph{regime} synthetic ranking $\rho{\approx}0.80$ are consistent with a failure pattern that is not tied to a single shift magnitude (computed by \texttt{scripts/analyze\_tta\_adaptation.py}). Recent TTA-pitfall work~\cite{zhao2023pitfalls,press2024rdumb} corroborates that TTA can fail systematically on non-vision domains.
